\chapter{Equalization}

\section{Overview}

This chapter introduces equalization techniques used to
compensate for channel loss and inter-symbol interference

\section{Rx CTLE Equalization}
Start from a typical CTLE response,
\begin{equation}
H_{\mathrm{CTLE}}(s)
=
A_0\frac{1+s/\omega_z}{1+s/\omega_p},
\qquad
\omega_z<\omega_p
\end{equation}

\begin{equation}
H_{\mathrm{CTLE}}(s)
=
A_0\frac{\omega_p}{\omega_z}
\left[
1-\frac{\omega_p-\omega_z}{s+\omega_p}
\right]
\end{equation}
The first term, \(1\), represents the signal itself, while the second term, \(1/(s+\omega_p)\), represents a low-pass filter. \\
Therefore, the CTLE can be interpreted as the original signal minus its low-pass-filtered component.\\
Now we take a look at how CTLE cancel post-cursor in time domain. First write the impulse response of CTLE, which is inverse-transform of transfer function \(H_{\mathrm{CTLE}}(s)\)

\begin{equation}
h_{\mathrm{CTLE}}(t)
=
K\delta(t)-Ce^{-\omega_p t}u(t)
\end{equation}

where

\begin{equation}
K=A_0\frac{\omega_p}{\omega_z},
\qquad
C=A_0\frac{\omega_p}{\omega_z}(\omega_p-\omega_z).
\end{equation}

For channel pulse response \(p_{\mathrm{ch}}(t)\),

\begin{equation}
p_{\mathrm{eq}}(t)
=
p_{\mathrm{ch}}(t)*h_{\mathrm{CTLE}}(t)
\end{equation}

thus

\begin{equation}
p_{\mathrm{eq}}(t)
=
Kp_{\mathrm{ch}}(t)
-
C\int_{-\infty}^{t}
p_{\mathrm{ch}}(\tau)
e^{-\omega_p(t-\tau)}\,d\tau.
\end{equation}

Define

\begin{equation}
m(t)
=
\int_{-\infty}^{t}
p_{\mathrm{ch}}(\tau)
e^{-\omega_p(t-\tau)}\,d\tau,
\end{equation}

then

\begin{equation}
p_{\mathrm{eq}}(t)
=
Kp_{\mathrm{ch}}(t)-Cm(t).
\end{equation}

Since \(m(t)\) is a low-pass filtered version of
\(p_{\mathrm{ch}}(t)\), the fast-changing main cursor is less
cancelled, while the slowly varying post-cursor tail is more strongly
cancelled.\\
From Chapter 2
\begin{equation}
p_{\mathrm{ch}}[n]
=
p_{\mathrm{ch}}(t_0+nT).
\end{equation}
Define

\begin{equation}
p_{\mathrm{eq}}[n]
=
p_{\mathrm{eq}}(t_0+nT).
\end{equation}

Since

\begin{equation}
p_{\mathrm{eq}}(t)
=
Kp_{\mathrm{ch}}(t)-Cm(t),
\end{equation}

we have

\begin{equation}
p_{\mathrm{eq}}[n]
=
Kp_{\mathrm{ch}}[n]
-
Cm(t_0+nT),
\end{equation}

and

\begin{equation}
p_{\mathrm{eq}}[0]
=
Kp_{\mathrm{ch}}[0]
-
Cm(t_0).
\end{equation}

To show

\begin{equation}
\frac{p_{\mathrm{eq}}[n]}{p_{\mathrm{eq}}[0]}
<
\frac{p_{\mathrm{ch}}[n]}{p_{\mathrm{ch}}[0]},
\end{equation}

substitute the above expressions:

\begin{equation}
\frac{Kp_{\mathrm{ch}}[n]-Cm(t_0+nT)}
{Kp_{\mathrm{ch}}[0]-Cm(t_0)}
<
\frac{p_{\mathrm{ch}}[n]}
{p_{\mathrm{ch}}[0]}.
\end{equation}

After cross multiplication and cancellation,

\begin{equation}
p_{\mathrm{ch}}[0]\,m(t_0+nT)
>
p_{\mathrm{ch}}[n]\,m(t_0).
\end{equation}

Equivalently,

\begin{equation}
\frac{m(t_0+nT)}{p_{\mathrm{ch}}[n]}
>
\frac{m(t_0)}{p_{\mathrm{ch}}[0]}.
\end{equation}

Since \(m(t)\) is a low-pass filtered version of
\(p_{\mathrm{ch}}(t)\), it follows the slowly varying post-cursor
more strongly than the fast-changing main cursor. Therefore,

\begin{equation}
\frac{p_{\mathrm{eq}}[n]}{p_{\mathrm{eq}}[0]}
<
\frac{p_{\mathrm{ch}}[n]}{p_{\mathrm{ch}}[0]},
\qquad n=1,2,3,\ldots
\end{equation}


